% GENERATED FILE. Edit canonical lessons, then run npm run generate:books.
% canonical-source-hash: fnv1a64:3e0cba12f67d187f
% canonical-lessons: ML-C256-vinodam, ML-C256-kalippattam, ML-C256-abhiprayam, ML-C256-asayam, ML-C256-upadesam

\chapter{Entertainment, Toy, Opinion, Idea, Advice}
\label{ch:ml-entertainment-toy-opinion-idea-advice}
\hlchaptermodality{256}

\begin{chapteropening}
\textbf{By the end of this chapter:} \emph{I can say entertainment, a toy, an opinion, an idea and advice.}

Complete the last lesson of chapter 256: I can say entertainment, a toy, an opinion, an idea and advice.
\end{chapteropening}

\section[\texorpdfstring{\ml{വിനോദം} (vinōdaṁ)}{വിനോദം (vinōdaṁ)}]{\ml{വിനോദം} (vinōdaṁ) — entertainment, a pastime}

\label{lesson:ML-C256-vinodam}



\begin{quote}
Before the new one: say the Malayalam for a camera, then the Malayalam for a printer.
\end{quote}

\subsection*{You'll want to know: \ml{വിനോദം}}
\textbf{\ml{വിനോദം}} — \emph{vinōdaṁ} — \textquotedblleft{}entertainment, a pastime\textquotedblright{}.

\begin{grammarlens}[title={What you've built}]
One more word for this chapter.
\end{grammarlens}

\subsection*{Your turn}
\begin{itemize}
\raggedright
  \item \emph{Say it:} \emph{vinōdaṁ}
  \item \emph{Say it:} \emph{vinōdaṁ}, once more
  \item \emph{Say it:} say \emph{vinōdaṁ}
  \item \emph{Recall:} say the Malayalam for a camera, then the Malayalam for a printer, then say \emph{vinōdaṁ} again
\end{itemize}

\begin{tcolorbox}[breakable,colback=teal!4,colframe=teal!35!black,title={Before you move on}]
What does \ml{വിനോദം} mean? (\textquotedblleft{}Entertainment, a pastime\textquotedblright{}.) Say it once more.
\end{tcolorbox}
\section[\texorpdfstring{\ml{കളിപ്പാട്ടം} (kaḷippāṭṭaṁ)}{കളിപ്പാട്ടം (kaḷippāṭṭaṁ)}]{\ml{കളിപ്പാട്ടം} (kaḷippāṭṭaṁ) — a toy}

\label{lesson:ML-C256-kalippattam}



\begin{quote}
Before the new one: say the Malayalam for a printer, then the Malayalam for entertainment.
\end{quote}

\subsection*{You'll want to know: \ml{കളിപ്പാട്ടം}}
\textbf{\ml{കളിപ്പാട്ടം}} — \emph{kaḷippāṭṭaṁ} — \textquotedblleft{}a toy\textquotedblright{}.

\begin{grammarlens}[title={What you've built}]
One more word for this chapter.
\end{grammarlens}

\subsection*{Your turn}
\begin{itemize}
\raggedright
  \item \emph{Say it:} \emph{kaḷippāṭṭaṁ}
  \item \emph{Say it:} \emph{kaḷippāṭṭaṁ}, once more
  \item \emph{Say it:} say \emph{kaḷippāṭṭaṁ}
  \item \emph{Recall:} say the Malayalam for a printer, then the Malayalam for entertainment, then say \emph{kaḷippāṭṭaṁ} again
\end{itemize}

\begin{tcolorbox}[breakable,colback=teal!4,colframe=teal!35!black,title={Before you move on}]
What does \ml{കളിപ്പാട്ടം} mean? (\textquotedblleft{}A toy\textquotedblright{}.) Say it once more.
\end{tcolorbox}
\section[\texorpdfstring{\ml{അഭിപ്രായം} (abhiprāyaṁ)}{അഭിപ്രായം (abhiprāyaṁ)}]{\ml{അഭിപ്രായം} (abhiprāyaṁ) — an opinion}

\label{lesson:ML-C256-abhiprayam}



\begin{quote}
Before the new one: say the Malayalam for entertainment, then the Malayalam for a toy.
\end{quote}

\subsection*{You'll want to know: \ml{അഭിപ്രായം}}
\textbf{\ml{അഭിപ്രായം}} — \emph{abhiprāyaṁ} — \textquotedblleft{}an opinion\textquotedblright{}.

\begin{grammarlens}[title={What you've built}]
One more word for this chapter.
\end{grammarlens}

\subsection*{Your turn}
\begin{itemize}
\raggedright
  \item \emph{Say it:} \emph{abhiprāyaṁ}
  \item \emph{Say it:} \emph{abhiprāyaṁ}, once more
  \item \emph{Say it:} say \emph{abhiprāyaṁ}
  \item \emph{Recall:} say the Malayalam for entertainment, then the Malayalam for a toy, then say \emph{abhiprāyaṁ} again
\end{itemize}

\begin{tcolorbox}[breakable,colback=teal!4,colframe=teal!35!black,title={Before you move on}]
What does \ml{അഭിപ്രായം} mean? (\textquotedblleft{}An opinion\textquotedblright{}.) Say it once more.
\end{tcolorbox}
\section[\texorpdfstring{\ml{ആശയം} (āśayaṁ)}{ആശയം (āśayaṁ)}]{\ml{ആശയം} (āśayaṁ) — an idea}

\label{lesson:ML-C256-asayam}



\begin{quote}
Before the new one: say the Malayalam for a toy, then the Malayalam for an opinion.
\end{quote}

\subsection*{You'll want to know: \ml{ആശയം}}
\textbf{\ml{ആശയം}} — \emph{āśayaṁ} — \textquotedblleft{}an idea\textquotedblright{}.

\begin{grammarlens}[title={What you've built}]
One more word for this chapter.
\end{grammarlens}

\subsection*{Your turn}
\begin{itemize}
\raggedright
  \item \emph{Say it:} \emph{āśayaṁ}
  \item \emph{Say it:} \emph{āśayaṁ}, once more
  \item \emph{Say it:} say \emph{āśayaṁ}
  \item \emph{Recall:} say the Malayalam for a toy, then the Malayalam for an opinion, then say \emph{āśayaṁ} again
\end{itemize}

\begin{tcolorbox}[breakable,colback=teal!4,colframe=teal!35!black,title={Before you move on}]
What does \ml{ആശയം} mean? (\textquotedblleft{}An idea\textquotedblright{}.) Say it once more.
\end{tcolorbox}
\section[\texorpdfstring{\ml{ഉപദേശം} (upadēśaṁ)}{ഉപദേശം (upadēśaṁ)}]{\ml{ഉപദേശം} (upadēśaṁ) — advice}

\label{lesson:ML-C256-upadesam}



\begin{quote}
Before the new one: say the Malayalam for an opinion, then the Malayalam for an idea.
\end{quote}

\subsection*{You'll want to know: \ml{ഉപദേശം}}
\textbf{\ml{ഉപദേശം}} — \emph{upadēśaṁ} — \textquotedblleft{}advice\textquotedblright{}.

\begin{grammarlens}[title={What you've built}]
One more word for this chapter.
\end{grammarlens}

\subsection*{Your turn}
\begin{itemize}
\raggedright
  \item \emph{Say it:} \emph{upadēśaṁ}
  \item \emph{Say it:} \emph{upadēśaṁ}, once more
  \item \emph{Say it:} say \emph{upadēśaṁ}
  \item \emph{Recall:} say the Malayalam for an opinion, then the Malayalam for an idea, then say \emph{upadēśaṁ} again
\end{itemize}

\begin{tcolorbox}[breakable,colback=teal!4,colframe=teal!35!black,title={Before you move on}]
What does \ml{ഉപദേശം} mean? (\textquotedblleft{}Advice\textquotedblright{}.) Say it once more.
\end{tcolorbox}
